\documentclass[a4paper,12pt]{article}
\usepackage[utf8]{inputenc}
\usepackage[spanish, mexico]{babel}
\usepackage{graphicx}
\usepackage{booktabs}
\usepackage[default]{lato}
\usepackage{titling}
\setlength{\droptitle}{-4cm}
\usepackage{datatool}	

\DTLsetseparator{,}
\DTLloadrawdb[keys={variable,valor}]{datos}{datos.csv}


\title{Reporte Diario de Obra}
\author{Departamento de Supervisión}
\date{\today}

\begin{document}
	
	\maketitle
	
	\section*{Datos Generales}
	\noindent \textbf{Obra:} \DTLfetch{datos}{variable}{obra}{valor} \\
	\textbf{Ubicación:} \DTLfetch{datos}{variable}{ubicacion}{valor} \\
	\textbf{Encargado:} \DTLfetch{datos}{variable}{encargado}{valor} \\
	\textbf{Fecha:} \today \\
	
	\section*{Condiciones del Día}
	\noindent \textbf{Clima:} \DTLfetch{datos}{variable}{clima}{valor} \\
	\textbf{Cantidad de Trabajadores:} \DTLfetch{datos}{variable}{trabajadores}{valor}
	
	\section*{Avance de Obra}
	El avance acumulado de la obra hasta la fecha es del \DTLfetch{datos}{variable}{avance}{valor}.  
	La actividad principal realizada hoy fue: \textit{\DTLfetch{datos}{variable}{actividadprincipal}{valor}}.
	
	\section*{Observaciones y Problemas}
	Durante la jornada se presentaron los siguientes inconvenientes:  
	\DTLfetch{datos}{variable}{problemas}{valor}  
	
	\section*{Acciones Tomadas}
	Para mitigar los problemas detectados, se han implementado las siguientes acciones:  
	\DTLfetch{datos}{variable}{acciones}{valor}  
	
	\section*{Conclusión}
	Se continuará con la planificación prevista, ajustando los tiempos según la disponibilidad de materiales.
	
\end{document}
